\clearpage
\section{Unrestricted observational context congruence}
\label{sec:copying-general-congruence}
The equal-input-depth restriction of the preceding section is now removed.
Every original annotated term context preserves the existing observational
equivalence, for arbitrary input terms and finite variable/atom environments,
provided execution respects recursive graph isomorphism and one-step padding.
These are explicit obligations for the concrete copying UAM, not project
axioms or consequences of determinism. The proposition's terms, the Figure~16
translation and the observation class are unchanged.

The key new structural result compares interpreting a context with an argument
at its original allowance and interpreting it with the padded argument. The
two resulting graphs are proved recursively isomorphic after padding to a
sufficient common allowance. Combining this with the universal addressed
filling law and the cofinal observation criterion yields unrestricted semantic
congruence. The final equivalence of the private-cell and constant-result
programs is still pending: congruence preserves an established equivalence;
it does not supply that equivalence or the template proofs.

\subsection{One-step padding through the literal graph constructors}
\termdef{copying-constructor-padding}

The uniform relabelling presentation of padding commutes with literal serial
assembly up to ordinary graph isomorphism. Keep each serial pattern vertex,
component tag, component vertex and component edge. The inverse maps retain
the same data, so both inverse identities follow by splitting on the left or
right tag. Internal labels are unchanged, and each edge has the same raised
component label. For sources and targets, map composition reduces the claim
to an individual component port. Internal, input and output ports have the
same attachment to the original two-hole serial pattern on both sides.
Thus entire ordered incidence lists are preserved.

The proof splits depth zero from positive depth to respect the existing
alphabet representation. Replacing the uniform relabelling presentation by
its proved equality with the original raising operation gives recursive
isomorphism between raising a serial composite and composing its raised
components.


\lean{Copying/PaddingAssembly}{9}{94}{PaddingAssembly.serial_graph}

\lean{Copying/PaddingAssembly}{95}{98}{PaddingAssembly.serial}

Uniform relabelling also commutes with ordered parallel assembly, in this
case as graph equality. The carriers, enumerations and labels coincide;
equality of source and target functions follows by inspecting each internal,
input or output port in the fixed boundary embedding. Replacing uniform
relabelling by the original raising operation gives the corresponding exact
parallel padding equality. Binary tensor composition is its concatenating
boundary instance. The Boolean family of raised components equals the family
of the two raised graphs by its two defining cases.


\lean{Copying/PaddingAssembly}{99}{107}{PaddingAssembly.parallel_graph}

\lean{Copying/PaddingAssembly}{108}{114}{PaddingAssembly.parallel}

\lean{Copying/PaddingAssembly}{115}{123}{PaddingAssembly.tensor}

Padding commutes with outer-interface equality casts by eliminating those
equalities. A primitive corolla padded once equals the same primitive corolla
at the next allowance, by splitting on the zero or positive alphabet case.
Distributor forests satisfy the analogous equality: their unchanged graph
is relabelled by the primitive-label embedding at the selected allowance.

The environment distributor is the ordered tensor of its term and atom
forests. Tensor padding and both forest equalities give its padding law.
The single-use distributor adds only the existing interface cast; commute
padding with that cast and apply the environment-distributor law. No
contraction, weakening edge or single-use distributor is discarded.


\lean{Copying/PaddingAssembly}{124}{128}{PaddingAssembly.castInterfaces}

\lean{Copying/PaddingAssembly}{129}{132}{PaddingAssembly.primitive}

\lean{Copying/PaddingAssembly}{133}{137}{PaddingAssembly.distributor}

\lean{Copying/PaddingAssembly}{138}{143}{PaddingAssembly.environmentDistributor}

\lean{Copying/PaddingAssembly}{144}{149}{PaddingAssembly.singleUse}

Raising a box preserves its bound-variable, free-variable and atom metadata,
and raises its body once. This is the original recursive definition of
raising. For a lambda box, commute raising with the existing body-interface
cast as well. The variable bound by the lambda remains in the original first
slot.

A translated pair uses the prescribed variable/atom middle-block exchange.
Commute padding with its outer cast and with parallel assembly, then identify
the two Boolean component families. The exchange itself is unchanged, so
the resulting equality retains the exact Figure~16 pair wiring.


\lean{Copying/PaddingAssembly}{150}{153}{PaddingAssembly.box}

\lean{Copying/PaddingAssembly}{154}{158}{PaddingAssembly.lambdaBox}

\lean{Copying/PaddingAssembly}{159}{166}{PaddingAssembly.pair}

For a unary clause, compose the serial padding isomorphism for the complete
clause with serial congruence applied to the inner operation-prefix padding
isomorphism. The primitive and single-use equalities identify the remaining
components with the original clause at the next allowance.

For abstraction, use the same two serial steps and then the primitive,
lambda-box and single-use equalities. Its host is raised at allowance $d+1$
while its body is raised at $d$, preserving their one-level offset.

Application and each binary operation use two serial steps around their
ordered pair. Apply the primitive, pair and environment-distributor equalities
to obtain their respective recursive isomorphisms. Each binary operation
parameter is arbitrary, retaining the full operation alphabet.


\lean{Copying/PaddingAssembly}{167}{174}{PaddingAssembly.unary}

\lean{Copying/PaddingAssembly}{175}{182}{PaddingAssembly.abstraction}

\lean{Copying/PaddingAssembly}{183}{190}{PaddingAssembly.application}

\lean{Copying/PaddingAssembly}{191}{199}{PaddingAssembly.binary}

\subsection{Padding to any specified adequate allowance}

Padding a graph of allowance $d$ to $e+1$, where $d\leq e$, equals raising
its padding to $e$ once. Write $e=d+n$ and use the exact padding-to-addition
identity twice. The final equality is the defining successor clause of
iterated padding.


\lean{Copying/PaddingConstructors}{7}{16}{PaddingConstructors.padTo_succ}

Each one-step constructor law extends to any adequate allowance by induction
on the allowance bound. At equality of the original and requested allowance,
identity padding gives reflexive isomorphism. At a successor, rewrite each
padding with the preceding successor identity, raise the induction
isomorphism, and compose with the corresponding one-step constructor law.
This proves the unary, application and binary laws separately.

For abstraction, a body padded from $d$ to $e$ is boxed at $e+1$; its original
abstraction is padded from $d+1$ to $e+1$. The same induction maintains this
shift. Existence of each isomorphism is proved before the Lean definition
selects it; no unproved structural assumption is introduced.


\lean{Copying/PaddingConstructors}{17}{29}{PaddingConstructors.unary_to}

\lean{Copying/PaddingConstructors}{30}{43}{PaddingConstructors.abstraction_to}

\lean{Copying/PaddingConstructors}{44}{57}{PaddingConstructors.application_to}

\lean{Copying/PaddingConstructors}{58}{72}{PaddingConstructors.binary_to}

\subsection{Structural comparison at a sufficient common allowance}
\termdef{copying-common-allowance-iso}

For two scoped graphs $G,H$, define $\mathsf{AgreeAt}(G,H,N)$ to consist of
proofs that both depths are at most $N$, together with a recursive isomorphism
between their original graphs padded to $N$. This is an auxiliary structural
predicate, not a quotient of graphs or a replacement observation relation.

Agreement at $N$ widens to any $M\geq N$: pad the isomorphism to $M$ and
collapse both intermediate paddings by exact padding transitivity. Agreement
is reflexive at every adequate allowance, symmetric by reversing the
isomorphism, and transitive at the same allowance by composing the two
isomorphisms. Different proofs of the middle graph's adequate bound carry
no additional graph data.


\lean{Copying/CommonAllowanceIso}{9}{15}{CommonAllowanceIso.AgreeAt}

\lean{Copying/CommonAllowanceIso}{16}{22}{CommonAllowanceIso.widen}

\lean{Copying/CommonAllowanceIso}{23}{25}{CommonAllowanceIso.refl}

\lean{Copying/CommonAllowanceIso}{26}{29}{CommonAllowanceIso.symm}

\lean{Copying/CommonAllowanceIso}{30}{35}{CommonAllowanceIso.trans}

Agreement is preserved by unary interpretation at the same allowance. Commute
padding with the left unary graph, apply unary congruence to the given padded
argument isomorphism, and compose with the inverse right unary padding law.
For abstraction, use allowance $N+1$ and the corresponding body-to-host
padding law; both adequacy bounds increase by one. This preserves the full
metadata and body interface of each box.


\lean{Copying/CommonAllowanceIso}{36}{41}{CommonAllowanceIso.unary}

\lean{Copying/CommonAllowanceIso}{42}{50}{CommonAllowanceIso.abstraction}

The scoped application interpreter first pads its two arguments to their
maximum depth. Padding its result further to an adequate $N$ agrees up to
recursive isomorphism with applying the literal constructor directly to
both arguments padded to $N$. Use the application padding law, then collapse
each argument's two successive paddings to $N$. The binary interpreter has
the same proof with its selected binary operation. Thus the original
intermediate maximum is retained in the construction and eliminated only
by a proved equality of padded graphs.


\lean{Copying/CommonAllowanceIso}{51}{58}{CommonAllowanceIso.application_to}

\lean{Copying/CommonAllowanceIso}{59}{68}{CommonAllowanceIso.binary_to}

Suppose the two first arguments agree at $N$ and the two second arguments
also agree at $N$. Both scoped applications have maximum depth at most $N$.
Use the preceding interpreter-padding isomorphism on the left, application
congruence of the two argument isomorphisms, and the inverse interpreter law
on the right. This proves agreement of applications at $N$. The analogous
three-step composite proves agreement for every binary operation.


\lean{Copying/CommonAllowanceIso}{69}{77}{CommonAllowanceIso.application}

\lean{Copying/CommonAllowanceIso}{78}{88}{CommonAllowanceIso.binary}

Define common-allowance isomorphism by the existence of some $N$ at which
the two scoped graphs agree. Reflexivity uses the graph's own depth; symmetry
retains the witness and reverses its agreement. For transitivity, widen the
two witnessing agreements to the maximum of their witnesses, then compose
there. A scoped isomorphism, which already includes equality of the depths,
supplies common-allowance isomorphism by padding it to the maximum of those
depths. All witnesses remain actual recursively isomorphic padded graphs.


\lean{Copying/CommonAllowanceIso}{89}{92}{CommonAllowanceIso.CommonIso}

\lean{Copying/CommonAllowanceIso}{93}{94}{CommonAllowanceIso.refl}

\lean{Copying/CommonAllowanceIso}{95}{97}{CommonAllowanceIso.symm}

\lean{Copying/CommonAllowanceIso}{98}{101}{CommonAllowanceIso.trans}

\lean{Copying/CommonAllowanceIso}{102}{105}{CommonAllowanceIso.ofScoped}

Unary and abstraction interpretation preserve common-allowance isomorphism
by using the respective fixed-allowance laws. For application and binary
interpretation, first widen both argument comparisons to the maximum of
their witnesses, then apply the fixed-allowance constructor theorem.

Induction on the complete original annotated context now proves that its
constructor interpretation preserves common-allowance isomorphism. The hole
uses the supplied relation; lambda and unary use their preservation laws.
At either application or binary position, the fixed sibling is related to
itself and the selected subcontext uses the induction hypothesis. Every
original context constructor is covered.

Finally, a scoped graph agrees with its own padding to any adequate allowance:
choose that allowance as witness, remove identity padding on the already
padded side, and use reflexive isomorphism. Combined with context preservation,
this is the missing unequal-depth structural comparison.


\lean{Copying/CommonAllowanceIso}{106}{109}{CommonAllowanceIso.unary}

\lean{Copying/CommonAllowanceIso}{110}{113}{CommonAllowanceIso.abstraction}

\lean{Copying/CommonAllowanceIso}{114}{118}{CommonAllowanceIso.application}

\lean{Copying/CommonAllowanceIso}{119}{124}{CommonAllowanceIso.binary}

\lean{Copying/CommonAllowanceIso}{125}{136}{CommonAllowanceIso.apply}

\lean{Copying/CommonAllowanceIso}{137}{143}{CommonAllowanceIso.padTo}

\subsection{The unrestricted semantic congruence theorem}
\termdef{copying-general-context-congruence}

Let execution respect recursive isomorphism and one-step padding, and assume
$t$ and $u$ are equivalent in the existing all-allowances observational sense.
For every original annotated context $C$, the theorem proves
\[
 \mathsf{ObservationalEquiv}(C[t],C[u]).
\]
There is no premise equating the minimal depths of $t$ and $u$.

\paragraph{A common input allowance and one actual hole.}
Choose $d$ to be the maximum of the two minimal term depths. Let $G,H$ be the
unchanged translated-program wrappers at $d$. Unrestricted filling supplies
one certified graph context $D$, of some host allowance $q$, with local
allowance $d+n$, whose actual filling with any graph padded by $n$ agrees
with its constructor interpretation. The same address works for both $G$
and $H$.

\paragraph{Correspondence with the original filled terms.}
Each original term translation is common-allowance isomorphic to its padding
to $d$. Context preservation transports these comparisons through $C$.
The exact constructor-interpretation identity identifies their unpadded
endpoints with the original translations of $C[t]$ and $C[u]$. Compose the
inverse comparisons with the addressed filling laws. This gives finite
witness allowances $N_t,N_u$ comparing each actual filling of $D$ with its
original filled-term translation. The two witness allowances may differ.

\paragraph{The semantic comparison at arbitrarily large outputs.}
Given any requested lower threshold $b$, choose
$N=\max(b,\max(N_t,N_u))$. Widen both structural comparisons to $N$. Their
adequacy proofs ensure $q\leq N$; write $N=q+m$.
The input programs of $D$ padded by this extra $m$ are the original
translations of $t$ and $u$ at allowance $(d+n)+m$, by exact program-padding
coherence. The assumed observational equivalence supplies their comparison.
Padded graph-context congruence yields equivalence of the two actual fillings
at host allowance $N$.

The two widened graph isomorphisms relate those fillings to the original
translated-program wrappers for $C[t]$ and $C[u]$ at $N$. Execution isomorphism
invariance turns them into observational equivalences; compose the inverse
left one, the contextual comparison, and the right one. This establishes
the cofinal comparison for every $b$. The previously proved cofinal criterion,
using the explicit execution padding law as well, yields the required
all-allowances conclusion.

No template robustness or determinism premise is used. The theorem is an
inference from input equivalence under stated representation laws. It does
not assert those laws for a copying execution or assert the input equivalence
of Proposition~12.2.


\lean{Copying/GeneralContextCongruence}{9}{57}{GeneralContextCongruence.observationalEquiv_context}

The structural comparison can also be stated directly: interpreting any
original context with a graph at its original allowance, or with that graph
padded to a larger adequate allowance, gives common-allowance isomorphic
results. Apply context preservation to the proved comparison of a graph with
its own padding. This theorem has no execution premise and changes no term,
constructor, boundary permutation or observation context.


\lean{Copying/GeneralContextCongruence}{58}{63}{GeneralContextCongruence.apply_padTo}

\subsection{Original-program regression and completion boundary}

For every original annotated context and every finite variable/atom scope,
compare its direct filled-term translation with its constructor interpretation
when the proposition's argument is supplied at allowance three. The private-cell
program has depth three; the constant-result program has depth two and is
padded once. Apply the preceding structural theorem to each unchanged term
translation and rewrite its unpadded interpretation with the exact filled-term
identity. This proves both correspondences independently. It does not compare
the private-cell program with the constant-result program, so it neither
assumes nor disguises the desired final equivalence.

The general context-filling and observational-congruence development of former
Appendix~B.1 item~3 is now complete. Its concrete execution hypotheses are
consolidated into the copying-machine task, where they remain unproved.
Appendix~B.1 therefore has four outstanding work packages: remaining translation
and invariant laws; the full copying machine and its representation laws;
the eleven template families; and semantic weakening, Figure~22 and the final
proposition. This bookkeeping does not declare the application complete.


\lean{Copying/GeneralCongruenceExamples}{11}{23}{GeneralCongruenceExamples.proposition12_2_context_representations}
